\chapter{The weak-C2 R9 contact corner: a full-parent coupled-circuit certificate}
\label{component:pb18dc55f4f2f}
\noindent\textbf{Component }\path{pb18dc55f4f2f}\par
{\small\noindent Source: \path{R16/inherited_r15/proofs/r9_contact/reciprocal_contact_closure.tex}\par}
{\small\noindent SHA256: \path{b18dc55f4f2fb27fe0c98a7478243f59e43e1b958a2b565dd8e423ff64bf9dc0}\par}
{\small\noindent Inventory use: RC1-RC13\par}
\medskip
\section{The weak-C2 R9 contact corner: a full-parent coupled-circuit certificate}

\paragraph{Scope.}
Retain the original R9 physical stationary source and all inherited
intersections. Thus $B,a,h,e$ are positive infinitesimals,
$r=1-e$, $t$ and $j$ are units with $t_0>1$, and $n\ge0$.
Consider arbitrary independent $u,v>0$ infinitesimal, with
$\omega=u/v>0$ unrestricted. The following is a certificate for the
entire original parent, not only its contact graph. Its bounded Gram
coefficients contain no $\omega$ or $1/\omega$, so it covers every
positive testing ratio at once. The original axes/sign identities and
R13 remote theorem then close all tests on R9.

All formulas below refer to $F/4$. Use the unchanged WC3, WC6,
WC13--WC17a coefficients, and put
\[
 x=aB=k^2,\quad y=tB=l^2,\quad c=\sqrt{a/t}=k/l,
 \quad U=u^2,\ V=v^2,\quad q=q_s+\alpha,
\]
\[
 \mu=m-a^2q,\quad \lambda=\frac{\mu}{ta^2B},\quad
 K_X'=K_X-a\alpha,\quad K_E'=K_E+a\alpha/t.
\]
Indeed $\lambda=f_B-\alpha/y$, where the R14 exact joint-source
divisions give $f_B\sim1/2$ and $\alpha/B^2\sim t_0/2$. Since
$\alpha/y=(B/t)(\alpha/B^2)$ is infinitesimal, $\lambda\sim1/2>0$.
The same exact reductions give
$c_F\sim1/(2t_0)$, $d_4\sim t_0/[2(t_0-1)]$ and $b_4$ bounded.
These reductions hold on the literal $n=0$ face, and no division by
$n$ is used. In particular all source denominators used below are positive; no
comparison of source and testing rates occurs.

\paragraph{An exact scalar completion.}
At $Z=1$ use
\[
 A=uv,\quad B_c=ck\,v^2/u,\quad C_c=k^2/c,
 \quad G=kv,\quad d=B_c-G,\quad J=A-C_c,
\]
\[
 F_v=V-y,\qquad R_D=v(U-x),\qquad H_u=(a/t)V/u-u.
\]
Thus $AB_cC_c=G^3$. Completing the entire $\alpha U$ term into
$qH_u^2$ changes $f_B$ to $\lambda=f_B-\alpha/y$ and $f_G$
to $f_G'=f_G-2\alpha/y$. The exact source contact balances imply
\[
 f_{A0}=\lambda-2d_4x-b_4y/t^2,\qquad
 f_{C0}=\lambda+b_4y/t^2,\qquad
 f_G'=3\lambda-d_4x+b_4y/t^2.
 \tag{RC1}
\]
These identities follow from the value and both first derivatives
of the original scalar at $(u,v)=(k,l)$, since the extracted
$H_u,F_v$ squares vanish there; they can also be checked by the
finite source formulas. Using $A^2=UV$, $G^2=xV$ and $C_c^2=xy$
now gives the exact scalar identity
\[
 \mathcal S=qH_u^2+nr^2u^2H_u^2+
 \left(c_F+\frac{b_4x}{t^2}\right)F_v^2
 +\lambda L_3+d_4R_D^2+\frac{b_4}{t^2}vR_DF_v,
 \tag{RC2}
\]
where $L_3=A^2+B_c^2+C_c^2-3G^2$.
For example the two nontrivial remainders are precisely
\[
 (U-2x)A^2+xG^2=V(U-x)^2,
\]
\[
 (V-y)A^2+y(C_c^2-G^2)
       =V(U-x)(V-y)+x(V-y)^2.
\]
This uses no limiting expansion or approximate scalar positivity.

\paragraph{The complete original hard block.}
Keep the full WC4 state frame
\[
 E=tuW-vZ,\qquad N_y=(Y-huW)/\sqrt h,\qquad
 \xi=X+ruZ.
\]
Work initially at $Z=1$. Complete the two angular squares as
\[
 Q=\sqrt q\left(H_u+\frac{a}{t\omega}E-a\xi\right)
   =\sqrt q\,(avW-uZ-a\xi),
\]
\[
 N=\sqrt n\{r uH_u+u\xi\}.
 \tag{RC3}
\]
The second formula is homogenized by replacing $uH_u$ by
$[(a/t)V-U]Z$. It includes $n=0$ without division.
Subtracting $Q^2+N^2$ from the entire parent leaves the singular
hard quadratic exactly as the single positive rank-one packet
\[
 \mu\left(\frac{E}{t\omega}-\xi\right)^2.
 \tag{RC4}
\]
All other hard coefficients form a bounded quadratic
$\mathcal H_0(E,N_y,\xi)$ with residue
\[
 \operatorname{diag}\left(\frac1{2(t_0-1)},\frac{t_0}2,
                                     \frac12\right).
 \tag{RC5}
\]
This assertion concerns the full parent. To check that no hidden
$1/\omega$ remains, its only such term before subtraction is the
$m$-event coefficient $(1+hU)/(t\omega)$.
After removing its leading $m/(t\omega)$ part, the $E^2$ error is
$m(2hV+h^2UV)/t^2$, the $E\xi$ error is $-2mh uv/t$,
and the $EN_y$ coefficient is bounded by a bounded source factor
times $m\sqrt hV$. All remaining event coefficients have no
negative power of $\omega$: $uW=(v+E)/t$ and
$Y=h(v+E)/t+\sqrt hN_y$ make this explicit. The $A$ event uses
$Ah$, so no inverse power of $h$ survives. Subtracting $Q^2,N^2$
changes its bounded entries only by infinitesimals after RC4 is
removed. This proves RC5 directly from WC3, without a mixed inverse.

\paragraph{Exact centering of every first normal.}
Set
\[
 a_E=2Ah\,h/t^2-2e_1/t,\quad b_E=2b_2/t^2,
 \quad c_E=2b_4/t^2.
\]
Let $a_Y,b_Y,c_Y$ be the three coefficients of
$u^3w,u^3w^3,u^5w^3$ in $\mathcal R_Y$ in WC17a.
All are bounded. Original contact of the complete normals gives
\[
 2K_E'+a_Ex+b_Ey+c_Exy+2\lambda x=0,
\]
\[
 2K_X'+x_3y+x_5xy-2t\lambda x=0,
 \qquad 2K_Y+a_Yx+b_Yy+c_Yxy=0.
 \tag{RC6}
\]
These are source identities. With $p=ck/\omega$ (possibly
unbounded), one has $pG=kB_c$ and hence
\[
 pB_c-kG=(p+k)d.
\]
The entire residual normal coefficients after RC3 are exactly
\[
 \mathcal R_E'=2\lambda(p+k)d
 +(a_E+c_EV)R_D+(b_E+c_Ex)vF_v,
 \tag{RC7}
\]
\[
 \mathcal R_X'=-2t\lambda ck\,d
 +(2t\lambda c+x_5uv)R_D-2t\lambda cuJ
                    +u(x_3+x_5x)F_v,
 \tag{RC8}
\]
\[
 \mathcal R_Y=(a_Y+c_YV)R_D+(b_Y+c_Yx)vF_v.
 \tag{RC9}
\]
The actual $N_y$ coefficient is $\sqrt h\mathcal R_Y$.
The identity underlying RC8 is particularly useful:
\[
                  xu-ckG=cR_D-cuJ.
\]
The identities $UV-xy=V(U-x)+x(V-y)$ and RC6 prove all three
formulas. Thus no normal is inferred merely from its vanishing.

\paragraph{Absorb the unbounded ratio without comparing it.}
The elementary circuit identity and its exact coupled completion are
\[
 L_3=J^2+(1+2k/p)d^2,
\]
\[
 \begin{split}
 \lambda L_3+\lambda p^2E^2+2\lambda(p+k)Ed
 ={}&\lambda J^2+\frac{2\lambda k}{p}(pE+d)^2\\
 &+\lambda\{(p-k)E+d\}^2-\lambda k^2E^2.
 \end{split}                                                    \tag{RC10}
\]
All radicals in its positive squares are separated: $k/p=\omega/c$.
No boundedness assumption on $p$ has been used. Put
$R=\sqrt\lambda\{(p-k)E+d\}$.
The singular cross term in RC4 must be kept with RC8. Together they
are
\[
 -2t\lambda ck\,\xi(pE+d)
 =-2t\sqrt\lambda ck\,\xi R-2t\lambda ck^2E\xi.
 \tag{RC11}
\]
Both displayed coefficients are infinitesimal. The remaining
$\mu\xi^2$ term is bounded and infinitesimal. This exact cancellation
is the step that permits arbitrary independent rates of $c$ and
$\omega$.

The two coupled roots simplify after homogenization to
\[
 pE+dZ=kv(tcW-Z),\qquad
 R=\sqrt\lambda\,tk(cv-u)W,
\]
\[
 \sqrt{2\lambda k/p}\,(pE+dZ)
       =\sqrt{2\lambda/c}\,k\sqrt{uv}\,(tcW-Z).
 \tag{RC13}
\]
In particular their final forms contain no testing-ratio inverse.
Retain $Q^2,N^2,(2\lambda k/p)(pE+d)^2$ as exact positive squares.
Every remaining term is a bounded Gram on
\[
              (E,R_D,N_y,\xi,F_v,J,R).
\]
Its only nonzero off-diagonal residue is between $E$ and $R_D$:
$a_E\to-1/(t_0-1)$. Their full residue block is
\[
 \frac1{2(t_0-1)}
 \begin{pmatrix}1&-1\\-1&t_0\end{pmatrix}\succ0,
 \qquad \det=\frac1{4(t_0-1)}.
 \tag{RC12}
\]
The other diagonal residues are $t_0/2,1/2,1/(2t_0),1/2,1$.
All other entries are infinitesimal: this is explicit in RC7--RC11,
using $\sqrt h,c,u,v,k$ infinitesimal and all printed source arrays
bounded. The subtraction $-\lambda k^2E^2$ in RC10 does not change
the positive residue. The bounded-Gram lemma therefore supplies
finite tensor SOS for the entire parent.

\paragraph{Restoration and all-tests join.}
Homogenize all scalar roots by $Z$. The expressions RC3 are already
state-linear, and every other root becomes a state-linear expression;
no $Z$ inverse remains. Hence $Z=0$ is included. All divisions are
pure source or pure testing. Although $p$ and $k/p$ may be unbounded,
they occur only inside the printed exact positive square or the root
$R$, never as an entry of the bounded Gram. No ratio hypothesis on
$u/v$ has been used. The exact original gauges restore both
angle signs; real and imaginary parts give complex states. The final roots in RC3 and RC13 and all other Gram roots have no
testing inverse, so substitution extends the identity to either literal
axis directly. The inherited R13 R9 remote theorem covers the complement
of both angles infinitesimal.
These are exhaustive pure-testing cases; no testing-ratio split remains.
Thus this coupled certificate closes the whole R9 source shape,
with every original source gap and literal $n=0$ retained. No other
weak-C2 source shape is claimed here.
